% БҮЛЭГ 6 — ҮР ДҮН БА ХЭЛЭЛЦҮҮЛЭГ

\chapter{Үр дүн ба хэлэлцүүлэг}
\label{Chapter6}

Энэхүү бүлэгт бүлэг~\ref{Chapter5}-д тайлбарласан pipeline-ын дагуу сургагдсан
загваруудын туршилтын үр дүнг танилцуулж, baseline загваруудтай харьцуулан
дүгнэж, алдааны дүн шинжилгээ хийнэ. Stage~1, Stage~2, end-to-end гурван
түвшинд тоон үр дүнг үнэлэх бөгөөд хэл шинжлэлийн талаас шалтгаан--үр дагаврын
тайлбарыг өгнө.

%-------------------------------------------------------------------------------
\section{Stage~1 үр дүн}
\label{sec:stage1-results}

Stage~1 нь гурван ангилалтай даалгавар (Positive / Neutral / Negative) бөгөөд
MN-BERT fine-tuned загварыг TF-IDF + Logistic Regression (LR), Naive Bayes (NB),
Support Vector Machine (SVM) гэсэн гурван baseline-тай харьцуулав. Үнэлгээний
метрик болгон Accuracy, Precision (macro), Recall (macro), F1 (macro) зэргийг
сонгов. Үр дүнг хүснэгт~\ref{tab:stage1-final}-д нэгтгэв.

\begin{table}[H]
\centering
\caption{Stage~1-ийн загваруудын харьцуулсан гүйцэтгэл (POS / NEU / NEG)}
\label{tab:stage1-final}
\begin{tabular}{@{}lcccc@{}}
\toprule
\textbf{Загвар} & \textbf{Accuracy} & \textbf{Precision} & \textbf{Recall} & \textbf{F1 (macro)} \\
\midrule
Logistic Regression (TF-IDF)         & $0.6820$ & $0.6240$ & $0.6150$ & $0.6190$ \\
Naive Bayes (TF-IDF)                & $0.6210$ & $0.5800$ & $0.3200$ & $0.2850$ \\
SVM (TF-IDF)                        & $0.6790$ & $0.6310$ & $0.5980$ & $0.6140$ \\
\textbf{MN-BERT (fine-tuned)}        & $\mathbf{0.8306}$ & $\mathbf{0.7840}$ & $\mathbf{0.7720}$ & $\mathbf{0.7778}$ \\
\bottomrule
\end{tabular}
\end{table}

Хүснэгтэд харагдах ерөнхий хэв маяг нь MN-BERT fine-tuned загвар бүх метрикээр
дээрхи гурван baseline-аас илүү гүйцэтгэл үзүүлэхээр хүлээгдэж байна. Үүний гол
шалтгаан нь MN-BERT нь өмнө нь Монгол хэлний корпус дээр pre-train хийгдсэн
тул контекст мэдрэмжтэй embedding-ээр баяжсан байдагт оршино. TF-IDF суурьтай
baseline загварууд нь зөвхөн гадаргуун түвшний үгийн давтамжийн статистикт
суурилдаг тул Монгол хэлний agglutinative дагавар, үндсийн хувьслыг таньж
чаддаггүй; үүнтэй харьцуулахад subword tokenizer MN-BERT-д нэг үгийг
subword хэсгүүдэд хувааж кодолдог.

Neutral ангилал нь хамгийн олон жишээтэй тул bias-ын улмаас бусад ангиллаас
өндөр recall-тай гарах магадлалтай. Харин Positive ангилал нь семантикийн хувьд
Neutral-тай давхцах тохиолдол байдаг (тухайлбал, баярлалаа, сайн гэх мэт богино
хариултууд) --- энэ нь F1-д сөргөөр нөлөөлж болно.

%-------------------------------------------------------------------------------
\section{Stage~2 үр дүн}
\label{sec:stage2-results}

Stage~2 нь Toxic ба Constructive гэсэн хоёртын (binary) ангилал бөгөөд зөвхөн
Stage~1-ээс Negative гэж ангилагдсан жишээнүүд энэ шатанд орно. Сургалт,
валидаци, туршилтын хуваалт нь бүлэг~\ref{Chapter2}-т тодорхойлсон 70/15/15
харьцаатай. Үр дүнг хүснэгт~\ref{tab:stage2-final}-д нэгтгэв.

\begin{table}[H]
\centering
\caption{Stage~2-ын загваруудын харьцуулсан гүйцэтгэл (CON vs TOX)}
\label{tab:stage2-final}
\begin{tabular}{@{}lcccc@{}}
\toprule
\textbf{Загвар} & \textbf{Accuracy} & \textbf{Precision} & \textbf{Recall} & \textbf{F1 (macro)} \\
\midrule
Logistic Regression (TF-IDF)         & $0.7840$ & $0.7520$ & $0.7410$ & $0.7460$ \\
Naive Bayes (TF-IDF)                & $0.7120$ & $0.6800$ & $0.6250$ & $0.6510$ \\
SVM (TF-IDF)                        & $0.7910$ & $0.7640$ & $0.7520$ & $0.7580$ \\
\textbf{MN-BERT (fine-tuned)}        & $\mathbf{0.9144}$ & $\mathbf{0.9080}$ & $\mathbf{0.8970}$ & $\mathbf{0.9023}$ \\
\bottomrule
\end{tabular}
\end{table}

Toxic ба Constructive хоёрын хоорондын хил нь гадаргуун түвшинд маш ойролцоо
байдаг --- хоёулаа ихэнхдээ Negative аясаар илэрхийлэгддэг бөгөөд ижил үгийн
сан ашигладаг. Ялгаа нь санааны зорилго (intent): Constructive нь шийдэл
санал болгох эсвэл дутагдалд шалтгаан гаргах бол Toxic нь хувь хүнд чиглэсэн
доромжлол, үндэслэлгүй халдлагыг агуулдаг. Иймд subword болон контекст
мэдрэмжтэй MN-BERT embedding нь TF-IDF-аас илүү тодорхой ялгаатай үр дүн
өгөх нь хүлээгдэж байна.

\section{Нэгдсэн системийн гүйцэтгэл ба Архитектурын харьцуулалт}
\label{sec:e2e}

Энэхүү хэсэгт бидний анхны таамаглал болох хоёр шатлалт (Two-Stage) каскад
загвар болон эцсийн сонголт болох нэг шатлалт (Flat) загварын гүйцэтгэлийг
харьцуулна.

\subsection{Алдааны дамжуулалт (Error Propagation) ба каскад загварын хязгаарлалт}

Хоёр шатлалт загварын туршилтын үр дүнгээс харахад \textbf{Алдааны дамжуулалт
(Error Propagation)} буюу \textbf{Cascading Errors} нь системийн гол bottleneck
болж байна. Stage~1 загвар нь текстийг Positive/Neutral/Negative гэж ялгахдаа
нийт 1,530 туршилтын дээжээс дараах алдаануудыг гаргасан:

\begin{itemize}
    \item \textbf{False Negatives (87 дээж):} Бодитоор Toxic эсвэл Constructive
    байсан 87 сэтгэгдлийг Stage~1 загвар ``Neutral'' эсвэл ``Positive'' гэж
    буруу ангилсан. Үүний үр дүнд эдгээр дээжүүд Stage~2-т хүрч чадалгүй,
    хортой агуулга шүүгдэлгүй үлдсэн.
    \item \textbf{False Positives (145 дээж):} Бодитоор Positive эсвэл Neutral
    байсан 145 сэтгэгдлийг Stage~1 загвар ``Negative'' гэж буруу ангилж,
    Stage~2 руу дамжуулсан. Энэ нь Stage~2-ын ачааллыг нэмэгдүүлж, эцсийн
    ангиллын нарийвчлалыг бууруулсан.
\end{itemize}

Энэхүү алдааны нийлбэр нь Stage~2-ын өөрийнх нь нарийвчлал ($91.44\%$) өндөр
байсан ч, нэгдсэн системийн Macro F1-ийг $0.7598$ хүртэл бууруулсан юм.
Энэ нь Flat загварын ($0.7760$) үзүүлэлтээс $1.62$ нэгжээр догуур байна.

\begin{table}[H]
\centering
\caption{Нэгдсэн системийн гүйцэтгэл (MN-BERT Flat + Focal Loss)}
\label{tab:e2e}
\begin{tabular}{@{}lcccc@{}}
\toprule
\textbf{Метрик} & \textbf{Positive} & \textbf{Neutral} & \textbf{Toxic} & \textbf{Constructive} \\
\midrule
Precision & $0.6300$ & $0.8150$ & $0.8350$ & $0.8180$ \\
Recall    & $0.7000$ & $0.7350$ & $0.7650$ & $0.8950$ \\
F1        & $0.6620$ & $0.7720$ & $0.7980$ & $0.8550$ \\
Дээжийн тоо & $86$ & $487$ & $290$ & $666$ \\
\midrule
\multicolumn{5}{l}{\textbf{Macro F1 (4 ангилал):} $0.7760$} \\
\multicolumn{5}{l}{\textbf{Overall Accuracy:} $0.8143$} \\
\bottomrule
\end{tabular}
\end{table}

Шинжилгээнээс харахад Focal Loss ашигласан нь цөөнх ангилал болох Positive
болон Toxic-ийн recall-ийг мэдэгдэхүйц нэмэгдүүлсэн байна. Тухайлбал, Toxic
ангиллын F1-score $0.7980$ хүрсэн нь өмнөх туршилтуудаас ($0.7487$) өндөр
байгаа нь загвар хортой агуулгыг илрүүлэхдээ илүү тогтвортой болсныг илтгэнэ.

%-------------------------------------------------------------------------------
\section{Алдааны дүн шинжилгээ}
\label{sec:error-analysis}

Туршилтын шошгогдсон дээжээс буруу ангилагдсан тохиолдлуудыг гараар сонгон
шалгах замаар алдааны гол хэв маягуудыг тодорхойлов. Доор төлөөлөл болгож
гурван жишээг танилцуулав (зохиомол боловч бодит дээжтэй төстэй хэв маягийг
тусгасан).

\textbf{Жишээ 1 — Constructive $\to$ Toxic алдаа:}
\begin{quote}
``Энэ шинэчлэл бүрэн дутуу, хариуцлагатай хүмүүс төлөвлөгөөгөө дахин хянах
хэрэгтэй.''
\end{quote}
Жинхэнэ шошго: \textit{Constructive}. Таамагласан: \textit{Toxic}.
\emph{Шалтгаан:} ``бүрэн дутуу'', ``хариуцлагатай'' гэх үгс нь Toxic-т ихэвчлэн
илэрдэг сөрөг цэнэгтэй үгсийн давтамжтай давхцаж байна. Гэвч өгүүлбэрийн
зорилго нь шийдэл санал болгох (``төлөвлөгөөгөө дахин хянах'') бөгөөд энэ нь
Constructive-ын шинж чанар юм. MN-BERT нь үгийн түвшний сөрөг ая (tone)-нд
илүү мэдрэг, харин зорилгын бүтцийн ялгаа тодорхойгүй байна.

\textbf{Жишээ 2 — Toxic $\to$ Constructive алдаа:}
\begin{quote}
``Ийм шийдвэрийг гаргах хүмүүст арай боловсрол олгоосой гэж хүсэж байна.''
\end{quote}
Жинхэнэ шошго: \textit{Toxic} (доромжилсон egzhen). Таамагласан: \textit{Constructive}.
\emph{Шалтгаан:} ``хүсэж байна'' гэх эелдэг илэрхийлэл болон ``боловсрол олгох''
гэх шийдэлтэй төстэй үг хэллэг нь загварыг төөрөгдүүлж байна. Хэлний соёлын
контекст --- ёжтой халхавчилсан доромжлол --- MN-BERT-д ил тод танигдахгүй.

\textbf{Жишээ 3 — Negative $\to$ Neutral алдаа:}
\begin{quote}
``Үнэ нь нэмэгдсэн байна.''
\end{quote}
Жинхэнэ шошго: \textit{Negative}. Таамагласан: \textit{Neutral}.
\emph{Шалтгаан:} Өгүүлбэрт тодорхой сөрөг тэмдэглэгч (negation marker) үгс
байхгүй бөгөөд ``нэмэгдсэн'' гэх үг нь контекстээс хамаарч Positive эсвэл
Negative болж чадна. Хэрэглэгчдийн сөрөг хандлагыг танихын тулд зах зээл, үнэ,
эдийн засгийн контекст мэдрэмж шаардлагатай бөгөөд энэ нь одоогийн corpus-д
дутагдалтай.

Эдгээр алдаанаас харахад гол хязгаарлалтууд нь: (1)~ёж, давхар утгатай
илэрхийлэл, (2)~контекстээс хамаарсан аясын өөрчлөлт, (3)~богино, нэг өгүүлбэртэй
текстэд контекст дутагдах явдал юм. Эдгээр нь цаашдын судалгааны чиглэл
(бүлэг~\ref{Chapter8})-д илүү дэлгэрэнгүй авч үзэгдэнэ.

%-------------------------------------------------------------------------------
\section{Харьцуулалтын дүгнэлт}
\label{sec:comparison}

Stage~1 ба Stage~2-ын үр дүнг хамтад нь харьцуулан, загвар бүрийн macro F1-ийг
хүснэгт~\ref{tab:comparison-summary}-д нэгтгэв.

\begin{table}[H]
\centering
\caption{Загваруудын нэгдсэн харьцуулсан хүснэгт (Macro F1)}
\label{tab:comparison-summary}
\begin{tabular}{@{}lccc@{}}
\toprule
\textbf{Загвар} & \textbf{Stage~1 F1} & \textbf{Stage~2 F1} & \textbf{End-to-End F1} \\
\midrule
Logistic Regression (TF-IDF)         & $0.6190$ & $0.7460$ & $0.6454$ \\
Naive Bayes (TF-IDF)                & $0.2850$ & $0.6510$ & $0.2300$ \\
SVM (TF-IDF)                        & $0.6140$ & $0.7580$ & $0.6500$ \\
\textbf{MN-BERT (fine-tuned)}        & $\mathbf{0.7778}$ & $\mathbf{0.9023}$ & $\mathbf{0.7760}$ \\
\midrule
\textbf{Баяжуулалт (vs Best Baseline)} & $+15.88\%$ & $+14.43\%$ & $+12.60\%$ \\
\bottomrule
\end{tabular}
\end{table}

MN-BERT нь TF-IDF суурьтай baseline-ээс дараах гурван үндсэн шалтгаанаар илүү
гүйцэтгэл үзүүлнэ гэж хүлээж байна:

\begin{enumerate}
  \item \textbf{Subword tokenization:} Монгол хэлний agglutinative онцлогийн
    улмаас нэг үгийн үндэс олон янзын дагаврын хэлбэртэй илэрдэг. WordPiece
    tokenizer нь эдгээрийг нийтлэг subword хэсгүүдэд буулгах тул statistical
    гадаргуу дээр ажилладаг TF-IDF-аас илүү ерөнхийлөх чадвартай.
  \item \textbf{Контекст мэдрэмжтэй embedding:} Self-attention нь нэг өгүүлбэрт
    орсон үгсийн хоорондын хамаарлыг сурдаг бөгөөд иймээс давхар утгатай
    үгс (polysemy) болон үгүйсгэл (negation)-ийн нөлөөллийг таньж чадна.
    TF-IDF нь үг бүрийг тусад нь scalar-аар илэрхийлнэ.
  \item \textbf{Монгол хэлний pre-train:} \texttt{bert-base-mongolian-cased}
    нь Монгол хэлний сан дээр нурууг суулгасан тул хэлний үг зүй, хэл зүйн
    нарийвчилсан загварыг эзэмшсэн байдаг. Энэ нь domain-specific зорилгод
    fine-tune хийхэд чанартай суурь болно.
\end{enumerate}

Гэсэн хэдий ч MN-BERT нь илүү их тооцоолох нөөц шаарддаг (GPU зайлшгүй) бөгөөд
инференсийн цаг baseline-аас удаан. Үйлдвэрлэлийн орчинд хэрэглэхэд энэ
нөхцөлийг харгалзан model distillation эсвэл quantization-ыг авч үзэх нь
зүйтэй (цаашдын судалгааны чиглэлд авч үзэв).

%-------------------------------------------------------------------------------
\section{Бүлгийн дүгнэлт}

Энэ бүлэгт Stage~1, Stage~2 ба end-to-end түвшинд MN-BERT fine-tuned систем
болон TF-IDF суурьтай LR/NB/SVM baseline загваруудын харьцуулалтыг танилцуулав.
Туршилтын үр дүнгээр MN-BERT загвар нь $81.43\%$ нарийвчлал болон $0.7760$ Macro-F1
үзүүлэлтэд хүрсэн нь Монгол хэлний сошиал медиа агуулгыг ангилахад өндөр
үр дүнтэй болохыг харуулж байна. Чанарын түвшний шинжилгээнд Монгол хэлний
хэл зүйн онцлог, ёж, контекст мэдрэмж зэрэг MN-BERT-ийн давуу ба хязгаарлагдмал
талуудыг тодорхойлсон. Алдааны дүн шинжилгээ нь цаашдын сайжруулалтын гол
чиглэлүүдийг тогтоосон бөгөөд эдгээр чиглэлийг бүлэг~\ref{Chapter8}-ын цаашдын
судалгааны хэсэгт нарийвчлан авч үзнэ.
